\chapter{تمهيد}

لهذا الكتاب غاية واحدة بسيطة: كتاب كيمياء واحد متماسك يغطي كل شيء من
السنة الأولى في المدرسة إلى نهاية المرحلة الجامعية الأولى، وقد كُتب أصله
بالإنجليزية لقرّاء في كل أنحاء العالم، ثم نُقل هنا إلى العربية.

الأسلوب موجز وصارم: دروس مبنية على تعريفات وأمثلة وقضايا وطرائق
ومخططات، مع استدلال دقيق كلما كان في متناول المستوى المطروح، تتبعها تمارين
منتقاة بعناية. أما النتائج التي تُذكر بلا اشتقاق كامل فيُشار إليها صراحةً على
أننا نقبلها. وتُجمع الحلول الكاملة لكل التمارين في آخر الكتاب.
والمعادلات الكيميائية موزونة دائمًا، وكل عدد يُذكر لمادة حقيقية
مأخوذ من جدول مرجعي.

يتألف المقرر من أربعة كتب. يغطي الأول السنوات المدرسية الاثنتي عشرة كلها
في مجلد واحد ولا يدرّس المفهوم نفسه مرتين أبدًا: فكل فكرة
تُدرَّس مرة واحدة، في السنة التي تناسبها أولًا، والفصل اللاحق الذي
يحتاج إليها يبدأ بتذكير قصير، \emph{ما تعرفه من قبل}، قبل أن يمضي
أبعد. والكتب الثلاثة الأخرى هي سنوات الجامعة الثلاث.
وفي الكتاب الأول كل سنة مدرسية جزء مقسوم إلى فصول؛ وكل كتاب
جامعي سنة واحدة مقسومة إلى فصول. وكلما كانت السنة أبكر، كان القارئ الذي كُتبت
له أصغر: رسوم توضيحية أكثر، وخطوات محلولة أكثر تفصيلًا، وتمارين
ألطف.

\section*{كيف يُقرأ هذا الكتاب}

العبارات مرقّمة داخل كل فصل ومميَّزة بالألوان: التعريفات
بالأزرق، والمبرهنات بالأحمر، والقضايا ونظيراتها بالبرتقالي،
والطرائق --- وهي وصفات قصيرة للمهام المعتادة --- بالأخضر. والتمارين
متدرجة من ($\star$) للتطبيقات المباشرة إلى ($\star\star\star$) للمسائل
الأصعب. وتضم الأطر الملوّنة ما ليس جزءًا من الاستدلال: تذكيرًا
بما درّسه فصل سابق، أو طريقة إجراء عملية في
المختبر، أو صفحة من التاريخ، أو مخاطر منتج.

\section*{السلامة}

تُتعلَّم الكيمياء بمواد حقيقية، وبعضها خطر.
والتجارب الموصوفة في هذا الكتاب تُجرى في مختبر مدرسي أو
جامعي، على يد معلم أو بحضوره، مع معدات الوقاية
التي يوفرها المختبر. ولا يُقصد بأي منها أن يُجرى في
البيت.

\section*{المساهمة}

مصدر هذا الكتاب موجود في مستودع git عام:
\begin{center}
\url{https://github.com/bvirrion/one-chemistry-book}
\end{center}
والمساهمات --- من فصول جديدة، وتصحيحات، وشروح أفضل،
وتمارين إضافية --- مرحَّب بها؛ انظر \texttt{CONTRIBUTING.md} في
جذر المستودع.
